在最一般的意义下，电流是大量带电粒子沿确定方向发生整体定向运动的宏观表现。电能的输运和使用
已经成为现代社会的基础；如果所有电力驱动的设备突然停止工作，生活与生产会立即受到全面影响。
本讲从电流的定义出发，讨论物质导电的经典图像、元件的伏安特性、欧姆定律（Ohm's law）以及一般直流网络的
基尔霍夫定律（Kirchhoff's circuit laws）。下一讲还将在这些基础上研究更复杂的电阻网络。

导体和绝缘体的关键差别，在于物质内部是否存在能够在外场作用下长距离迁移的带电粒子，即载流子。
金属中的载流子主要是电子；电解液中正、负离子都可参与导电；半导体中则可能同时存在电子和空穴。
载流子的浓度、质量、电荷量和碰撞性质，以及温度、外电场、光照等外部条件，共同决定材料的导电性。
必须强调：金属导电的精确理论属于量子物理。下面的经典模型不能从第一性原理给出全部材料参数，
但能把宏观规律与载流子的运动联系起来，并说明经典图像在哪些条件下有效、又会在何处失效。

\subsection*{电流强度和电流密度}

电流并不只存在于接有电源的闭合回路中；只要带电粒子发生有方向的输运，就形成电流。取一个有向截面，
若在时间间隔 $\Delta t$ 内按规定正方向通过它的净电量为 $\Delta Q$，平均电流强度定义为
\begin{equation}
  \boxed{I=\frac{\Delta Q}{\Delta t}}.
\end{equation}
当时间间隔趋于零，瞬时电流是通过电量 $Q(t)$ 对时间的导数：
\begin{equation}
  \boxed{I(t)=\lim_{\Delta t\to0}\frac{\Delta Q}{\Delta t}
  =\frac{\mathrm dQ}{\mathrm dt}}.
\end{equation}
一般情况下 $I$ 随时间变化；若它不随时间变化，称为稳恒电流。反过来，已知电流随时间的关系后，
从 $t_1$ 到 $t_2$ 通过截面的净电量为
\begin{equation}
  \boxed{Q(t_1\to t_2)=\int_{t_1}^{t_2}I(t)\,\mathrm dt}.
\end{equation}
对于恒定电流，上式退化为 $Q=I(t_2-t_1)$。电流的国际单位是安培，
$1\,\mathrm A=1\,\mathrm C\,\mathrm s^{-1}$。

设导线横截面积为 $A$，单位体积内有 $n$ 个同种载流子，每个载流子的电量为 $q$，沿导线的平均定向
漂移速度为 $v_d$。在 $\Delta t$ 内能够穿过截面的载流子原来位于长为 $v_d\Delta t$ 的柱体中，
其数目为 $nAv_d\Delta t$，因而通过的电量是 $nqAv_d\Delta t$。代入定义得到
\begin{equation}
  \boxed{I=nqAv_d}.
\end{equation}
式中 $q$ 是带符号的电荷量，$v_d$ 是沿所选截面法向的速度分量。对电子，电荷与沿常规电流方向的
速度分量均为负，乘积为正。若只讨论大小，则写 $|I|=n|q|A|v_d|$，不可只对电荷取绝对值。

真实导线中的漂移速度通常很小。以横截面积 $A=1\,\mathrm{mm^2}$ 的纯银导线为例，若每个银原子提供
一个导电电子，则 $n\simeq6\times10^{28}\,\mathrm{m^{-3}}$。通有 $1\,\mathrm A$ 电流时
\begin{equation}
  \boxed{v_d=\frac{I}{neA}\approx1.0\times10^{-4}\,\mathrm{m\,s^{-1}}}.
\end{equation}
这个速度比金属中电子本身的无规则运动速度小许多个数量级。因此“电子像小球一样沿导线迅速跑到用电器”
并不是正确图像：电路中建立电场的过程很快，但单个电子的净漂移极慢。
金属电子的无规则运动速率不能用经典能量均分定理准确求得，必须由量子统计决定；其典型高速电子可达到
约光速的百分之一，而上面算出的净漂移仍只有 $10^{-4}\,\mathrm{m\,s^{-1}}$ 量级。

若截面较大，不同位置的载流子漂移可能并不相同，单一的 $I$ 不能描述局部流动。定义电流密度矢量
\begin{equation}
  \boxed{\boldsymbol j=nq\boldsymbol v_d}.
\end{equation}
它的方向按正电荷运动方向规定。对于电子，$q=-e$，所以 $\boldsymbol j$ 与电子漂移速度方向相反；
若有多种载流子同时运动，总电流密度是各类载流子贡献之和。穿过任意有向曲面 $S$ 的电流为
\[
  I=\iint_S\boldsymbol j\cdot\mathrm d\boldsymbol S.
\]
这正像流体的体积流量等于速度场在截面上的通量，只不过电流计算的是电荷通量。
例如河水中央常比近岸处流得快，必须先描述各点流速才能求总流量；非均匀导体中的电流也完全类似。

\subsection*{元件的伏安特性}

电路由元件按一定方式连接而成。对电路设计者而言，常可暂时不追究元件内部结构，而只研究端口处可测量
的电压和电流。最简单的是二端元件。图中先人为规定端口电压 $V$ 的正负方向和电流 $I$ 的测试方向；
这种规定不等于元件本身一定具有固定正、负极，计算出的负号只说明真实方向与规定方向相反。

\ivCharacteristicFigure
\diagramcaption{图 5.1：二端元件及其 $I$--$V$ 特性。标出的正、负号与箭头只是统一的测试约定。}

若元件稳定工作时，同一个端电压总对应同一个端电流，就可用函数 $I=I(V)$ 描述它的外部性能，称为
伏安关系；把它画在 $I$--$V$ 平面上，便得到伏安特性曲线。不同元件的曲线可以完全不同，线性电阻
只是其中最简单的一类。使用这种二端描述还隐含两个条件：一是所研究的回路区域内没有随时间变化的磁通，
否则端电压会受感应电场影响而与路径有关；二是元件内部不能持续积累净电荷，否则流入与流出的电流并不
相等。满足这些条件后，元件的伏安关系加上连接规律就足以决定整个稳恒电路的行为。


% EM teaching enrichment 20260927
\paragraph{工作点与两种电阻}
电动势 $\mathcal E$、内阻 $r$ 的电源驱动非线性元件时，电源端口满足 $U=\mathcal E-rI$。它与元件曲线 $I=I(U)$ 的交点才是实际工作点，不能把电源电动势直接当作元件端电压。
在该点，$R_{\rm dc}=U/I$ 是静态电阻，$r_d=\dd U/\dd I$ 是小变化对应的微分电阻。若元件满足 $U=aI+bI^3$，则两者分别为 $a+bI^2$ 与 $a+3bI^2$；它们仅在线性情形相同。微分关系只能描述工作点附近的小扰动。

\subsection*{Ohm 定律}

Ohm 定律给出一类导体两端电压与其中电流的实验关系。它并不直接解释导电的微观机制，却在广泛的材料
和工作条件下极其准确，因而成为电路分析中使用最频繁的规律之一。Georg Simon Ohm 在十九世纪系统
研究了这一关系；后来以他的姓氏命名电阻单位欧姆，并用纪念邮票表现电阻和公式。

\ohmHistoryFigure

在温度和其他外部条件不变时，若导体两端电压为 $U$、通过电流为 $I$，Ohm 定律写成
\begin{equation}
  \boxed{U=IR,\qquad I=\frac UR}.
\end{equation}
比例系数 $R$ 是电阻，单位为欧姆，$1\,\Omega=1\,\mathrm V/\mathrm A$。对长度为 $l$、横截面积为
$A$ 的均匀圆柱形导体，实验进一步给出
\begin{equation}
  \boxed{R=\rho\frac lA},
\end{equation}
其中 $\rho$ 是材料在给定状态下的电阻率。长度越长，载流子沿途经历的散射越多；截面积越大，可并行
输运电荷的通道越多，因而分别出现 $R\propto l$ 与 $R\propto1/A$。

\cylindricalResistorFigure
\diagramcaption{图 5.4：长度为 $l$、截面积为 $A$ 的均匀圆柱形电阻。电流方向与端电压测试方向一致。}

理想电阻的伏安关系是一条过原点的直线。若纵轴为 $I$、横轴为 $V$，直线斜率为 $1/R$；曲线越陡，
相同电压产生的电流越大，电阻越小。

\ohmicIVFigure
\diagramcaption{图 5.5：理想电阻的 $I$--$V$ 特性。直线斜率是电导 $1/R$。}

把宏观 Ohm 定律写成局部形式。对图 5.4 的均匀导体，内部电场 $E=U/l$，电流密度 $j=I/A$。
由 $U=IR$ 与 $R=\rho l/A$ 得
\begin{equation}
  \boxed{\boldsymbol j=\frac1\rho\boldsymbol E=\sigma\boldsymbol E},
\end{equation}
其中 $\sigma=1/\rho$ 为电导率。它不含导体的长度和截面积，直接描述材料局部的导电能力。上式只在
各向同性、线性且工作状态不改变的 Ohmic 材料中成立；对各向异性材料，$\sigma$ 一般应写成张量。


% EM teaching enrichment 20260927
\paragraph{非均匀导体的电阻积分与功率检验}
缓慢变截面的细导体可近似看成许多串接薄片，故
$\dd R=\rho(x)\dd x/A(x)$，$R=\int\rho(x)\dd x/A(x)$。若长度 $l$ 内截面积由 $A_0$ 线性增至 $2A_0$、电阻率不变，得到
\[
 R=\frac{\rho l}{A_0}\int_0^1\frac{\dd u}{1+u}
 =\frac{\rho l}{A_0}\ln2.
\]
这里采用一维电流近似；骤变截面附近的三维电流拥挤不能忽略时，应另求电流场。
局部耗散功率密度为 $\boldsymbol j\cdot\boldsymbol E=j^2/\sigma$，对体积积分应回到 $UI=I^2R$。单位和总功率是检查非均匀导体积分的有效方法。

\subsection*{导电的微观机制——简单模型}

先构造一个纯粹唯象的“黏滞阻力”模型。设载流子定向运动时受到与漂移速度成正比、方向相反的平均阻力
$\boldsymbol F_r=-\alpha\boldsymbol v_d$。恒定电场中达到稳定漂移后，电场力与阻力平衡：
\[
  q\boldsymbol E-\alpha\boldsymbol v_d=0,
  \qquad \boldsymbol v_d=\frac q\alpha\boldsymbol E.
\]
于是
\[
  \boldsymbol j=nq\boldsymbol v_d=\frac{nq^2}{\alpha}\boldsymbol E.
\]
与微观 Ohm 定律比较可得
\begin{equation}
  \boxed{\sigma=\frac{nq^2}{\alpha}}.
\end{equation}
这个模型把极复杂的散射过程压缩成一个平均阻力系数。它虽然不是真实金属中电子运动的精确图像，却能
说明为什么稳定状态下 $\boldsymbol j$ 与 $\boldsymbol E$ 同向且成正比。

它还提示了暂态过程：电压突然改变时，载流子具有惯性，漂移速度不能无延迟地跳变。若保留质量 $m$，
运动方程为 $m\dot{\boldsymbol v}_d=q\boldsymbol E-\alpha\boldsymbol v_d$，电流会在特征时间
$\tau=m/\alpha$ 内趋近新的稳态值。通常金属中的 $\tau$ 极短，所以普通低频电路把这一过程近似为
瞬时完成；例题 5.2 将把它定量写出。


% EM teaching enrichment 20260927 supplementary

\paragraph{从建立过程看稳定电流}
这个模型也给出电场接通后的响应。初始漂移速度为零时，
\[
 m\frac{\dd\vect v_d}{\dd t}=q\vect E-\alpha\vect v_d,
 \qquad \vect v_d(t)=\frac{q\vect E}{\alpha}(1-\mathrm e^{-t/\tau}),\quad \tau=\frac m\alpha.
\]
乘以 $nq$ 后得到 $\vect j(t)=\sigma\vect E(1-\mathrm e^{-t/\tau})$，其中 $\sigma=nq^2/\alpha$。
达到 $t=\tau$ 时电流约为最终值的 $63\%$，约数个 $\tau$ 后才可用稳态欧姆定律。
负载流子的漂移方向与电场相反，但 $q\vect v_d$ 与电场同向，故电导率仍为正。
这一时间尺度描述载流子的局部松弛，不能把它当作信号沿整根导线传播的时间。

\subsection*{电子的扩散和漂移}

金属中的电子并非只做缓慢定向运动。它们本来就在晶体中高速无规则运动，外电场只在这种运动上叠加
一个很小的平均漂移。电子在两次散射之间近似自由运动，受电场加速；散射会把先前积累的定向速度随机化，
大量电子的统计平均才形成稳定电流。

\drudeFigure
\diagramcaption{图 5.6：导体内电子的无规则运动与定向漂移。蓝色曲线表示一条示意轨迹，金色圆点表示晶格散射中心。}

设自由电子数密度为 $n$、电子电荷量的绝对值为 $e$、质量为 $m$，两次有效散射之间的平均时间为
$\tau$。在弱电场中，电子由电场获得的平均漂移速度大小约为 $v_d=eE\tau/m$。因此
\begin{equation}
  \boxed{\boldsymbol j=\sigma\boldsymbol E,\qquad
  \sigma=\frac{ne^2\tau}{m}}.
\end{equation}
这就是经典 Drude 模型给出的电导率。若从一次碰撞后零定向速度开始计算到下一次碰撞，简单的匀加速
平均会出现 $1/2$；但真实碰撞时间服从分布，Boltzmann 方程的弛豫时间近似给出上式的标准系数。

该模型也帮助理解金属电阻率随温度升高而增大的现象。金属中载流子浓度 $n$ 大致不变；升温增强晶格
振动，使散射更频繁、$\tau$ 变小，于是 $\sigma$ 减小而 $\rho$ 增大。半导体则不同：升温会显著增加
载流子浓度，$n$ 增大的效果常超过 $\tau$ 减小的效果，因此其电阻通常随温度升高而减小。
在较粗略的经典估计中，也可暂把平均自由程 $\ell$ 看作不变：温度升高使电子典型速率增大，
$\tau\sim\ell/\bar v$ 因而减小，仍得到同一趋势；但金属中更准确的解释需要量子统计与晶格散射理论。

推导上述电导率公式时隐含了几项条件：电子在一次自由飞行中由外场获得的定向速度远小于其无规则运动速度；
散射机制和平均碰撞时间不被外场显著改变；散射足够频繁；载流子浓度不随外场改变。任何一项遭到破坏，
伏安关系都可能偏离直线。稀薄气体在强电场中的放电就是典型例子：电子在两次碰撞之间获得很大能量，
可激发或电离气体分子，使载流子浓度发生剧烈变化，依次出现暗放电、辉光放电和电弧区。

\gasDischargeFigure
\diagramcaption{图 5.7：气体放电的典型电压--电流曲线。横轴为对数电流标度，各区段明显不满足线性 Ohm 定律。}

经典模型仍不能精确描述金属电子，但能把实测电导率换算成有意义的微观尺度。以银为例，
$\sigma\simeq6.3\times10^7\,\mathrm{S\,m^{-1}}$，取
$n\simeq6.0\times10^{28}\,\mathrm{m^{-3}}$，由 $\sigma=ne^2\tau/m$ 得到
\begin{equation}
  \boxed{\tau=\frac{m\sigma}{ne^2}\approx3.8\times10^{-14}\,\mathrm s}.
\end{equation}
若电子典型无规则运动速率取 $\bar v\sim10^6\,\mathrm{m\,s^{-1}}$，平均自由程约为
\begin{equation}
  \boxed{\ell\sim\bar v\tau\approx3.8\times10^{-8}\,\mathrm m}.
\end{equation}
它远大于晶格中原子间的典型距离，说明电子并非每经过一个原子就发生一次经典小球式碰撞；这一事实正是
后来建立能带理论和量子输运理论的重要线索。例如取原子间距约 $0.3\,\mathrm{nm}$，
上述 $38\,\mathrm{nm}$ 自由程约为其 $127$ 倍。这与朴素的“逐原子碰撞”图像相差很大。


% EM teaching enrichment 20260927
\paragraph{浓度不均匀时还有扩散流}
无规则热运动在均匀密度下彼此抵消；若载流子数密度 $n$ 随位置变化，粒子通量可近似写成
$\boldsymbol\Gamma_{\rm diff}=-D\nabla n$，相应电流为 $q\boldsymbol\Gamma_{\rm diff}$。加上漂移项有
\[
 \boldsymbol j=nq\boldsymbol v_d-qD\nabla n.
\]
这里 $D$ 是扩散系数，单位 $\mathrm{m^2/s}$。对电子 $q=-e$，扩散电流方向与电子扩散方向相反。
稳态零电流可以来自漂移与扩散相互抵消，而非各项分别为零。因此含浓度梯度的区域不能只据 $\boldsymbol j=0$ 就断言 $\boldsymbol E=0$；前面的纯欧姆关系省略了这一驱动力。

\subsection*{简单电阻网络}

电阻最基本的连接方式是串联和并联。串联支路中各电阻通过同一电流 $I$，总电压是各段电压之和；
并联支路中各电阻接在同一对节点之间，端电压相同，总电流是各支路电流之和。

\seriesParallelResistorsFigure
\diagramcaption{图 5.8：电阻的串联与并联。判断连接方式的依据是节点关系，而不是图形看起来是否横排或竖排。}

对 $n$ 个串联电阻，$U_i=IR_i$，所以
\[
  U=U_1+\cdots+U_n=I(R_1+\cdots+R_n).
\]
等效电阻为
\begin{equation}
  \boxed{R_{\rm s}=R_1+R_2+\cdots+R_n}.
\end{equation}
对 $n$ 个并联电阻，公共端电压为 $U$，第 $i$ 条支路电流为
\begin{equation}
  \boxed{I_i=\frac{U}{R_i}}.
\end{equation}
总电流
$I=\sum_i I_i=U\sum_i(1/R_i)$，故等效电阻满足
\begin{equation}
  \boxed{\frac1{R_{\rm p}}=\frac1{R_1}+\frac1{R_2}+\cdots+\frac1{R_n}}.
\end{equation}
串联等效电阻大于任一分电阻，并联等效电阻小于任一分电阻。若网络能够逐级辨认出纯串联或纯并联，
即可反复用上述串并联公式化简；但含桥路、多个电源或交叉耦合支路的网络通常不能这样处理。


% EM teaching enrichment 20260927
\paragraph{分压、分流及仪表加载}
两电阻串联时 $U_2=UR_2/(R_1+R_2)$；两电阻并联、总电流为 $I$ 时 $I_1=IR_2/(R_1+R_2)$。分流分子是另一支路的电阻，来源是先求共同电压。
若用内阻 $R_V$ 的电压表测 $R_2$，被测支路已变为 $R_2\parallel R_V$，读数是
\[
 U_V=U\frac{R_2R_V/(R_2+R_V)}{R_1+R_2R_V/(R_2+R_V)}.
\]
只有 $R_V\gg R_2$ 且加载误差可忽略时才恢复理想分压。先把仪表画成实际元件，再讨论它的极限，可防止“测量不影响电路”的未经检验假设。

\subsection*{电阻网络定理与基本术语}

一般电路是由元件按拓扑关系连接成的网络。图中的每个方框都代表一个元件，它可以是电阻、电源或更复杂
的二端器件。两个或更多元件的连接点称为节点；同一段理想导线上的各点电势相同，属于同一个节点。
相邻两节点之间包含元件的一段称为支路。若一个节点只连接两条支路，通常只是串接中的平凡节点；三条及
以上支路相接时会发生电流分流。沿若干支路从一个节点出发，最后回到原节点的闭合路径称为回路。

\networkTermsFigure
\diagramcaption{图 5.9：电路分析中的元件、节点、支路和回路。虚线只标示所选回路，不是导线。}


% EM teaching enrichment 20260927
\paragraph{画图之前先识别节点}
所有用无电阻导线直接连通的点标为同一节点；交叉但未连接的导线不合并。随后把每个元件记成“节点 $a$ 到节点 $b$ 的支路”，图形折叠或拉伸便不会改变电路。
含 $N$ 个节点、$b$ 条支路的连通网络中，选 $N-1$ 条无闭环支路可连接全部节点，称为生成树。每补入一条其余支路，恰增加一个独立回路，因此独立回路数为 $b-N+1$。这是方程数量的来源，也能用来检查是否把交叉线误当节点。

\subsection*{Kirchhoff 电流定律（KCL）}

Kirchhoff 电流定律是电荷守恒在集中参数电路中的表达：任一节点流入电流之和等于流出电流之和；
等价地，按统一符号约定计入后，节点各支路电流的代数和为零。未知支路电流的方向可以预先任意假设，
若解得负值，说明真实方向与假设方向相反。

\kclKvlFigure
\diagramcaption{图 5.10：左图为 KCL 的支路电流约定，右图为同一网络的支路电压和回路方向。所有字母标注均不覆盖元件。}

按左图箭头，对四个节点分别可写为
\[
\begin{array}{ll}
N_1:&0=-i_1-i_2-i_3,\\
N_2:&0=i_1+i_4-i_6,\\
N_3:&0=i_2-i_4-i_5,\\
N_4:&0=i_3+i_5+i_6.
\end{array}
\]
最后一式中的 $i_6$ 由底边支路方向确定。把四式相加得到恒等式 $0=0$，因此含 $N$ 个节点的连通网络
只有 $N-1$ 个独立 KCL 方程。少写会丢失约束，多写则只是重复已有信息。


% EM teaching enrichment 20260927
\paragraph{节点暂时储存电荷时的完整守恒式}
用闭面包围一个小节点，流入减流出等于 $\dd Q_{\rm node}/\dd t$。稳恒状态右边为零，才得到常用 KCL。
在集中参数模型中，若节点对地有不可忽略电容 $C$，就应把 $C\dd V/\dd t$ 作为一条支路电流列入。例如通过电阻 $R$ 给该节点充电时，
\[
 \frac{\mathcal E-V}{R}=C\frac{\dd V}{\dd t}.
\]
常用节点法不是否定电荷积累，而是把积累明确归入电容元件。先选电流正向再写等式，解出的负数不应被擅自改成正数。

\subsection*{Kirchhoff 电压定律（KVL）}

Kirchhoff 电压定律是稳恒电路能量守恒的等价表达：沿任一闭合回路，各支路电压的代数和为零。使用时
先规定每个支路电压的正方向，再选定绕行方向；沿支路从正端走向负端记为一个符号，反向走则记为相反符号。
对图 5.10 右侧标出的四个回路，可写成
\[
\begin{array}{ll}
L_1:&0=-v_1+v_2+v_4,\\
L_2:&0=-v_2+v_3-v_5,\\
L_3:&0=-v_4+v_5-v_6,\\
L_4:&0=v_1+v_6-v_3.
\end{array}
\]
四式相加同样得到恒等式，所以它们并不独立。前三个小回路已经覆盖全部支路，外围回路 $L_4$ 可由前三式
线性组合得到。对连通网络，独立回路数等于“支路数 $-$ 节点数 $+1$”。KCL、KVL 再配合每个元件的
伏安关系，原则上可以求出全部支路电流和元件电压。


% EM teaching enrichment 20260927
\paragraph{一圈电压与一圈能量如何对应}
对电动势 $\mathcal E$、内阻 $r$、外电阻 $R$ 的闭合回路，沿电流绕行，源内非静电作用提升每单位电荷的能量 $\mathcal E$，电阻段合计降低 $I(R+r)$，故
$\mathcal E-Ir-IR=0$。
乘以 $I$ 得电源输出功率 $\mathcal EI=I^2r+I^2R$，同时核对了符号与能量。
若回路有变化磁通，应加入感应电动势；不能一面把感应效应遗漏，一面仍以纯静电势差的闭合和为零处理全部电场。

\subsection*{节点电压法}

直接把所有支路电流和元件电压都列为未知量，哪怕小网络也会产生很大的方程组。节点电压法先选一个节点
作为参考地，把其电势定为零；随后只把其余独立节点电势作为未知量。支路电压由两端节点电势之差自动给出，
因此 KVL 自动满足；再对未知节点写 KCL，就得到数量最少的一组方程。

\nodeMethodFigure
\diagramcaption{图 5.11：节点电压法示例。取 $c$ 为参考地，理想电压源固定 $a$ 点电势，只剩 $b$ 点电势未知。}

在图中令 $\varphi_c=0$，电压源使 $\varphi_a=E_0$，设 $\varphi_b=e$。以流入 $b$ 点为正：
从 $a$ 经 $R_1$ 流入的电流是 $(E_0-e)/R_1$，电流源注入 $I_0$，从 $b$ 经 $R_2$ 流出的电流是
$e/R_2$。KCL 给出
\[
  \frac{E_0-e}{R_1}+I_0-\frac e{R_2}=0.
\]
解得
\[
  \boxed{e=\frac{E_0/R_1+I_0}{1/R_1+1/R_2}}.
\]
节点电势一旦求出，每个元件两端的电压和各支路电流便可由其伏安关系依次得到。参考节点原则上可以任意
选择；经验上，选与较多理想电压源相连的节点作参考地，往往能减少辅助未知量，使代数最简。

% EM teaching enrichment 20260927
\paragraph{悬浮电压源与解后的核对}
若理想电压源位于两个未知节点 $a,b$ 之间，不能把其电流写成“电压除以零电阻”。可把两节点合成超节点，对流向外界的支路写 KCL，再补 $V_a-V_b=\mathcal E$。
例如 $a,b$ 分别经 $R_a,R_b$ 接地，则
$V_a/R_a+V_b/R_b=0$，与电压约束联立得
$V_a=\mathcal E R_a/(R_a+R_b)$、$V_b=-\mathcal E R_b/(R_a+R_b)$。
负节点电势是相对参考点而言，并非异常。解后应检查各节点电流守恒，以及全部元件吸收功率的代数和为零。

